\section{Discussão e alcance dos resultados}

A análise permitiu avaliar duas questões complementares. A primeira foi se as características disponíveis fornecem informação para estimar a popularidade observada. A redução do erro em relação às referências constantes indica que sim, dentro da coorte analisada. A segunda foi se os descritores de letras acrescentam informação ao Baseline. A comparação pareada mostrou melhora pequena, mas consistente, especialmente ao examinar a regressão aditiva com Ridge.

O resultado central pode ser sintetizado da seguinte forma: em 58.633 composições operacionais do desenvolvimento, avaliadas por validação cruzada agrupada e aninhada, o bloco Tradicional aditivo alcançou MAE de 6,8906, RMSE de 8,9016 e $R^2$ de 0,5612. Frente ao Baseline da mesma família, houve melhora das três métricas, redução do P95 do erro e diminuição agregada dos erros superiores a 20 pontos. A melhora do MAE foi inferior a um décimo de ponto e, por isso, deve ser descrita como contribuição incremental modesta.

A comparação não demonstra que propriedades das letras causem maior popularidade. Tampouco identifica uma propriedade textual específica como responsável pelo ganho: as 196 variáveis foram avaliadas como bloco, e muitas são correlacionadas. A investigação de efeitos individuais exigiria procedimentos próprios e cuidados adicionais com seleção de variáveis, dependência e interpretação de coeficientes.

A escassez de composições com resposta elevada é uma limitação relevante. Embora a amostra total seja numerosa, apenas 149 composições tinham popularidade de pelo menos 70, das quais 11 estavam entre 80 e 89. Esse desequilíbrio explica por que não se deve confundir precisão média com desempenho para grandes valores da resposta. A análise não permite decompor quanto do erro nessa região resulta de escassez amostral, atributos insuficientes, forma do modelo ou características não observadas.

Há outras limitações. A seleção por disponibilidade de letra modifica o perfil da base; a composição é definida operacionalmente; a auditoria de hashes exclui sobreposição textual exata entre splits, mas não substitui revisão humana do pareamento ou detecção de paráfrases e versões próximas. Além disso, os descritores LFM2b e os textos LRCLIB podem não representar exatamente a mesma versão da letra. Esses pontos restringem a generalização e devem permanecer explícitos na monografia.

As métricas não devem ser comparadas diretamente às de rodadas anteriores com menor amostra externa de validação, pois os denominadores e a composição dos conjuntos mudaram. Nesta rodada, as comparações internas são alinhadas: os blocos usam as mesmas observações e partições. O holdout originalmente avaliado na V0.6 não é uma nova confirmação independente, e seus alvos não foram carregados nesta etapa.

Do ponto de vista computacional, a execução completa durou aproximadamente 23 minutos e salvou 50 modelos de folds. As previsões foram reproduzidas com diferença numérica máxima igual a zero. Os arquivos de entrada permaneceram inalterados, e nenhuma configuração interna falhou. Foram registrados avisos de categorias não observadas no treino, tratados pela regra de codificação prevista. Essa verificação demonstra reprodutibilidade computacional dos artefatos; não comprova, por si só, ausência de todos os possíveis vieses científicos.

Por fim, esta etapa não utilizou embeddings na coorte ampliada. O bloco Tradicional já incorpora descritores de letras, mas a contribuição de representações produzidas por uma rede pré-treinada continua sendo uma pergunta a investigar. Uma etapa posterior deverá preservar a coorte e as partições, acrescentar essas representações e comparar os blocos Profundo e Combinado com as referências aqui estabelecidas. Até lá, não há evidência nesta rodada de que embeddings melhorem ou piorem o desempenho.

O principal avanço é, portanto, uma referência estatística reproduzível e interpretável: sabe-se quantas composições foram avaliadas, como os modelos foram selecionados, em quanto os descritores melhoraram as métricas e onde permaneceram as maiores dificuldades. Esses resultados delimitam o que foi alcançado e orientam a próxima comparação, sem antecipar suas conclusões.
